\noindent\SourceLeaf{12}
ex art. 15 Theoriae Motus Corporum Coelestium patet, si perihelium incremento exiguo $\mathrm d\Pi$, atque angulus $\varphi$ incremento $\mathrm d\varphi$ correcti concipiantur, valores correctos anomaliarum mediarum fieri
\[
\begin{gathered}
M+m\mathrm d\Pi+n\mathrm d\varphi\\
M'+m'\mathrm d\Pi+n'\mathrm d\varphi\\
M''+m''\mathrm d\Pi+n''\mathrm d\varphi\\
M'''+m'''\mathrm d\Pi+n'''\mathrm d\varphi
\end{gathered}
\]
siquidem potestates productaque correctionum $\mathrm d\Pi$, $\mathrm d\varphi$ negligere liceat. Habebuntur itaque tres expressiones motus medii diurni respectu perihelii, puta
\[
\begin{gathered}
\frac{M'-M}{t'-t}+\frac{m'-m}{t'-t}\mathrm d\Pi+\frac{n'-n}{t'-t}\mathrm d\varphi\\[5pt]
\frac{M''-M'}{t''-t'}+\frac{m''-m'}{t''-t'}\mathrm d\Pi+\frac{n''-n'}{t''-t'}\mathrm d\varphi\\[5pt]
\frac{M'''-M''}{t'''-t''}+\frac{m'''-m''}{t'''-t''}\mathrm d\Pi+\frac{n'''-n''}{t'''-t''}\mathrm d\varphi
\end{gathered}
\]
e quarum aequalitate incognitae $\mathrm d\Pi$, $\mathrm d\varphi$ eruentur. Quibus in expressione una substitutis, prodibit motus medius correctus, unde per methodum notam etiam semiaxis maior determinabitur. Denique substitutis valoribus correctionum $\mathrm d\Pi$, $\mathrm d\varphi$ in expressione pro aliqua anomalia media, huius valor correctus, atque inde epocha longitudinis mediae at tempus arbitrarium transferenda sponte demanabunt.

Ceterum facile patet, quo pacto etiam perturbationum tum periodicarum tum saecularium, si opus videatur, ratio haberi possit. Scilicet ab illis tantummodo longitudines in orbita datas, antequam in calculum introducantur, purgare, per posteriores autem longitudinem perihelii atque excentricitatem pro epocha arbitraria suppositam ad singula tempora $t$, $t'$, $t''$, $t'''$ transferre oportebit.
\Article{6}
Iam videamus, quomodo haec ad determinationem orbitae e quatuor oppositionibus observatis adhiberi possint. Sit longitudo planetae in ecliptica oppositione prima $\alpha$, latitudo geocentrica $\beta$, distantia terrae a Sole $R$, longitudo nodi ascendentis approximata $\AscendingNode$, inclinatio plani orbitae ad eclipticam approximata $i$. Quantitates ad oppositiones reliquas spectantes per characteres similes indice uno, duobus vel tribus distinctos denoto. Ex $\alpha$, $\AscendingNode$, $i$ deducetur per for%
\SourceLeaf{13}mulam notam longitudo in orbita, similisque calculus tribus longitudinibus reliquis in ecliptica superstruetur. Ex his quatuor longitudinibus in orbita per praecepta art. praec. derivabuntur elementa elliptica, atque hinc per formulas notas quatuor radii vectores $r$, $r'$, $r''$, $r'''$. Statuendo iam
\[
\frac{R\sin\beta}{r}=\sin(\beta-\gamma),\qquad \frac{R'\sin\beta'}{r'}=\sin(\beta'-\gamma')\quad\text{etc.}
\]
erunt manifesto $\gamma$, $\gamma'$, $\gamma''$, $\gamma'''$ quatuor longitudines heliocentricae, e latitudinibus geocentricis deductae. Derivemus easdem e longitudinibus, ponendo
\[
\operatorname{tang}i\sin(\alpha-\AscendingNode)=\operatorname{tang}\delta,\qquad \operatorname{tang}i\sin(\alpha'-\AscendingNode)=\operatorname{tang}\delta'\ \text{etc.}
\]
patetque, esse debere $\gamma=\delta$, $\gamma'=\delta'$, $\gamma''=\delta''$, $\gamma'''=\delta'''$, si pro $\AscendingNode$ atque $i$ valores veri accepti essent, siquidem omnibus quatuor oppositionibus (quae problema plusquam determinatum reddunt) per eandem ellipsin satisfacere possibile fuerit. Quod si secus eveniat, correctiones $\mathrm d\AscendingNode$, $\mathrm di$ quantitatibus $\AscendingNode$, $i$ applicandae sequenti modo elicientur. Quatenus hae correctiones tamquam quantitates exiguae primi ordinis spectantur, quarum potestates productaque negligere licet, manifesto adhibitis valoribus $\AscendingNode+\mathrm d\AscendingNode$, $i+\mathrm di$, loco latitudinum $\gamma$, $\delta$, $\gamma'$, $\delta'$ etc. prodibunt aliae huius formae
\[
\begin{array}{ll}
\gamma+a\mathrm d\AscendingNode+b\mathrm di,&\delta+c\mathrm d\AscendingNode+f\mathrm di\\
\gamma'+a'\mathrm d\AscendingNode+b'\mathrm di,&\delta'+c'\mathrm d\AscendingNode+f'\mathrm di\quad\text{etc.}
\end{array}
\]
ubi coëfficientes $c$, $f$, $c'$, $f'$ etc. invenientur per formulas
\[
\begin{array}{ll}
c=-\frac12\sin2\delta\operatorname{cotang}(\alpha-\AscendingNode),&f=\dfrac{\sin2\delta}{\sin2i}\\[5pt]
c'=-\frac12\sin2\delta'\operatorname{cotang}(\alpha'-\AscendingNode),&f'=\dfrac{\sin2\delta'}{\sin2i}\quad\text{etc.}
\end{array}
\]
Coëfficientes $a$, $b$ etc. quoque per operationes analyticas determinare possemus, sed praefero methodum sequentem. Calculum modo indicatum denuo repeto, adhibendo loco longitudinis nodi $\AscendingNode$ aliam exigua quantitate ad lubitum electa ab illa discrepantem, unde valores numerici coëfficientium $a$, $a'$ etc. facile eruentur; nova eiusdem calculi repetitio, accepta inclinatione paullulum mutata simili modo valores coëfficientium $b$, $b'$ etc. suppeditabit. Simul hoc modo patebit, quantas correctiones elementa elliptica passura sint, a correctionibus exiguis quantitatum $\AscendingNode$, $i$.

\SourceLeaf{14}
His ita factis, pro determinandis correctionibus $\mathrm d\AscendingNode$, $\mathrm di$ \textit{quatuor} adsunt aequationes
\[
\begin{gathered}
\gamma-\delta+(a-c)\mathrm d\AscendingNode+(b-f)\mathrm di=0\\
\gamma'-\delta'+(a'-c')\mathrm d\AscendingNode+(b'-f')\mathrm di=0\\
\gamma''-\delta''+(a''-c'')\mathrm d\AscendingNode+(b''-f'')\mathrm di=0\\
\gamma'''-\delta'''+(a'''-c''')\mathrm d\AscendingNode+(b'''-f''')\mathrm di=0
\end{gathered}
\]
quibus exacte satisfieri vix poterit propter observationum errores inevitabiles; quapropter valores incognitarum $\mathrm d\AscendingNode$, $\mathrm di$ maxime idonei per principia in Sect. III. Libri secundi Theoriae Motus Corporum Coelestium explicata determinabuntur. Hinc simul, per ea, quae modo diximus, correctiones reliquorum elementorum facile derivabuntur,
\Article{7}
Liceat hisce praeceptis adhuc quasdam observationes adiicere.

I. Vix opus erit admonitione, formulas nostras ei quoque casui inservire posse, ubi nodi non spectantur tamquam quiescentes. si modo valor longitudinis nodi suppositus ab epocha sua ad quatuor tempora rite reducatur, atque hi quatuor valores diversi pro singulis quatuor locis in formulis istis adhibeantur. Idem valet de inclinatione, siquidem ipsius variatio saecularis innotuerit, ipsiusque rationem habere operae pretium videatur. Si insuper perturbationes latitudinis heliocentricae in calculum introducere placet, hae in aequationibus praecedentibus manifesto vel a $\gamma$, $\gamma'$, $\gamma''$, $\gamma'''$ subduci, vel ipsis $\delta$, $\delta'$, $\delta''$, $\delta'''$ adiici debebunt.

III. Quoties inclinatio orbitae atque excentricitas modicae sunt, coëfficientes $a$, $b$, $a'$, $b'$ etc. tam parvi evadunt, ut ipsos negligere, adeoque computo secundo atque tertio supra praescripto supersedere liceat. Tunc inventis correctionibus $\mathrm d\AscendingNode$, $\mathrm di$ (nisi forte perparvae evaserint) calculum longitudinum in orbita novum valoribus correctis longitudinis nodi atque inclinationis, puta $\AscendingNode+\mathrm d\AscendingNode$, $i+\mathrm di$, superstruere, atque calculum elementorum ellipticorum ad normam art. 5 repetere oportebit, siquidem longitudines in orbita correctae a prioribus non correctis pluribus minutis secundis diversae prodierint. Ceterum vix umquam opus erit, coëfficientes $m$, $n$, $m'$ etc. denuo computare, quippe quorum valores iam per hypothesin primam satis exacti inveniuntur.
\SourceLeaf{15}
\Article{8}
Ad maiorem illustrationem praeceptorum praecedentium, calculum integrum, per quem systema elementorum tertium ex oppositionibus annorum 1805, 1807, 1808, 1809 determinatum est, hic apponam. Supposui longitudinem nodi ascendentis pro initio anni 1803, $172^0\,28'\,46''8=\AscendingNode$, inclinationem orbitae $34^0\,37'\,31''5$; reduxi illam ad tempora singularum oppositionum, addendo praecessionem $2'\,26''01$, $3'\,37''50$, $4'\,39''12$, $5'\,37''11$. Hinc derivavi ex longitudinibus latitudines heliocentricas
\[
\begin{array}{rrr}
\delta=-33^0&40'&50''63\\
\delta'=+28&14&51.24\\
\delta''=+27&20&55.86\\
\delta'''=-4&52&28.44
\end{array}
\]
\[
\begin{array}{ll}
c=+0.1252,&f=-0.9870\\
c'=-0.3366.&f'=+0.8917\\
c''=+0.3609.&f''=+0.8727\\
c'''=+0.6803,&f'''=-0{,}1811
\end{array}
\]
Porro prodeunt argumenta latitudinis
\[
\begin{array}{rrr}
257^0&25'&6''73\\
56&24&5.96\\
126&2&55.07\\
188&36&2.69
\end{array}
\]
e quibus elementa elliptica elicere oportet. Ad hunc finem statuo angulum $\varphi=14^0\,10'\,4''08$ (uti in systemate elementorum secundo inventus erat), longitudinemque perihelii, pro initio anni 1803, $=121^0\,5'\,22''1$. Perihelium (perinde ut nodus) respectu stellarum fixarum immobile supponitur, adeoque ipsius distantia a nodo ascendente constans $=308^0\,36'\,35''3$. Hinc habemus per formulas notas:
\[
\begin{array}{rrr@{\quad}rrr}
v=308^0&48'&31''43,&M=328^0&15'&45''08\\
v'=107&47&30.66,&M'=79&46&27.05\\
v''=177&26&19.77,&M''=175&54&28.87\\
v'''=239&59&27.39,&M'''=266&29&57.59
\end{array}
\]
\SourceLeaf{16}
\[
\begin{array}{ll}
m=-0.68517,&n=+1.18569\\
m'=-1.06482,&n=-2.01318\\
m''=-1.59701,&n=-0.12921\\
m'''=-1.18352,&n=+1.93947
\end{array}
\]
Denotando igitur correctionem longitudinis perihelii per $\mathrm d\Pi$, atque correctionem anguli $\varphi$ per $\mathrm d\varphi$, erit motus medius sidereus

ab oppositione prima usque ad secundam intra dies $521.140706$
\[
111^0\,30'\,41''97-0.37965\mathrm d\Pi-3.19887\mathrm d\varphi
\]
ab oppositione secunda usque ad tertiam intra dies $449.277674$
\[
96^0\,8'\,1''82-0.53219\mathrm d\Pi+1.88397\mathrm d\varphi
\]
ab oppositione tertia usque ad quartam intra dies $422.786667$
\[
90^0\,35'\,28''72+0.41349\mathrm d\Pi+2.06868\mathrm d\varphi
\]
Hinc eruuntur tres expressiones pro valore medio diurno sidereo
\[
\begin{gathered}
770''31398-0.0007285\mathrm d\Pi-0.0061382\mathrm d\varphi\\
770.30718-0.0011845\mathrm d\Pi+0.0041933\mathrm d\varphi\\
771.37892+0.0009780\mathrm d\Pi+0.0048930\mathrm d\varphi
\end{gathered}
\]
e quarum comparatione demanant aequationes:
\[
\begin{gathered}
0=0''00680+0.0004560\mathrm d\Pi-0.0103315\mathrm d\varphi\\
0=-1.07174-0.0021625\mathrm d\Pi-0.0006997\mathrm d\varphi
\end{gathered}
\]
atque hinc $\mathrm d\Pi=-488''82$, $\mathrm d\varphi=-20''92$, ac proin
\begin{center}\fontsize{9}{11}\selectfont\setlength{\tabcolsep}{3pt}\begin{tabular}{lr}
longitudo perihelii correcta pro initio anni 1803,&$120^0\,57'\,13''28$,\\
valorque correctus anguli $\ .\ .\ .\ .\ .\ .$&$\varphi=14^0\,9'\,43''16$\\
Porro fit motus medius diurnus sidereus $\ .\ .\ .\ .\ .$&$770''7985$\\
anomalia media correcta pro oppositione prima $\ .$&$328^0\,20'\,55''20$\\
longitudo media pro eadem epocha $\ .\ .\ .\ .\ .$&$89^0\,20'\,34''49$
\end{tabular}\end{center}
Denique derivatur ex motu medio sidereo logarithmus semiaxis maioris $0.4420439$, adeoque logarithmus semiparametri $0.4152361$. Iam quum habeamus valores correctos anomaliarum verarum
\SourceLeaf{17}
\[
\begin{array}{rrr}
v=308^0&56'&40''25\\
v'=107&55&39.48\\
v''=177&34&28.59\\
v'''=240&7&36.21
\end{array}
\]
computantur logarithmi radiorum vectorum
\[
\begin{aligned}
\log r&=0.3531088\\
\log r'&=0.4492406\\
\log r''&=0.5369700\\
\log r'''&=0.4716739
\end{aligned}
\]
E tabulis Solaribus porro habemus
\[
\begin{aligned}
\log R&=9.9937332\\
\log R'&=0.0039862\\
\log R''&=0.0065917\\
\log R'''&=0.0011160
\end{aligned}
\]
unde tandem deducimus
\[
\begin{array}{rrr}
\gamma=-33^0&39'&48''15\\
\gamma'=+28&15&0.73\\
\gamma''=+27&20&9.07\\
\gamma'''=-4&52&53.99
\end{array}
\]

Ut eruantur coëfficientes $a$, $b$ etc., formo \textit{hypothesin secundam}, retinendo inclinationem sed augendo longitudinem nodi uno minuto primo, ut sit pro initio anni 1803, $172^0\,29'\,46''8$. Statuendo dein, ut in hypothesi prima, longitudinem perihelii pro eadem epocha $=121^0\,5'\,22''1$, inveniuntur anomaliae verae
\[
\begin{array}{rrr}
v=308^0&48'&40''93\\
v'=107&47&34.06\\
v''=177&26&22.25\\
v'''=239&59&15.00
\end{array}
\]
atque hinc (statuendo quoque $\varphi=14^0\,10'\,4''08$) anomaliae mediae
\SourceLeaf{18}
\[
\begin{array}{rrr}
M=328^0&15'&51''59\\
M'=79&46&30.67\\
M''=175&54&32.83\\
M'''=266&29&42.93
\end{array}
\]
(Ceterum ad hunc computum non opus est, methodum vulgarem adhibere, sed sufficit, valoribus anomaliarum mediarum in hypothesi prima inventis addere producta e coëfficientibus $m$, $m'$, $m''$, $m'''$ positive sumtis in differentias inter valores respectivos anomaliarum verarum in utraque hypothesi, puta
\[
\begin{aligned}
+9''50\times0.68517&=+6''51\\
+3.40\times1.06482&=+3.62\\
+2.48\times1.59701&=+3.96\\
-12.39\times1.18352&=-14.66)
\end{aligned}
\]
Hinc retinendo valores coëfficientium $m$, $n$ etc. eruitur
\begin{center}\fontsize{9}{11}\selectfont\setlength{\tabcolsep}{3pt}\begin{tabular}{p{.23\linewidth}p{.70\linewidth}}
$\mathrm d\Pi=-468''21$,&longitudo perihelii initio anni 1803, $120^0\,57'\,33''89$\\
$\mathrm d\varphi=-20''62$,&$\varphi=14^0\,9'\,43''46$
\end{tabular}\end{center}
\begin{center}\fontsize{9}{11}\selectfont\setlength{\tabcolsep}{3pt}\begin{tabular}{ll}
motus medius diurnus sidereus&$=770''7761$\\
logarithmus semiaxis maioris&$=0.4420523$\\
longitudo media in oppositione prima&$=89^0\,20'\,47''85$
\end{tabular}\end{center}
\[
\begin{array}{rrr}
\gamma=-33^0&39'&51''10\\
\gamma'=+28&15&1.27\\
\gamma''=+27&20&10.97\\
\gamma'''=-4&52&54.66
\end{array}
\]

Tandem formo \textit{hypothesin tertiam} statuendo $\AscendingNode=172^0\,28'\,46''6$, $i=34^0\,38'\,31''5$, unde perinde ut ante deducitur
\begin{center}\fontsize{9}{11}\selectfont\setlength{\tabcolsep}{3pt}\begin{tabular}{lr}
longitudo perihelii pro initio anni 1803,&$120^0\,55'\,34''46$\\
&$\varphi=14^0\,9'\,52''63$\\
motus medius diurnus sidereus&$=770''8398$\\
logarithmus semiaxis maioris&$=0.4420283$\\
longitudo media in oppositione prima&$=89^0\,20'\,20''65$
\end{tabular}\end{center}
\SourceLeaf{19}
\[
\begin{array}{rrr}
\gamma=-33^0&39'&35''63\\
\gamma'=+28&15&5.20\\
\gamma''=+27&20&9.32\\
\gamma'''=-4&52&52.65
\end{array}
\]

Comparatio valorum latitudinum heliocentricarum $\gamma$, $\gamma'$, $\gamma''$, $\gamma'''$ in tribus hypothesibus inventarum producit
\[
\begin{array}{ll}
a=-0.0492,&b=+0.2087\\
a'=+0.0212,&b'=+0.0745\\
a''=+0.0317,&b''=+0.0042\\
a'''=-0.0112,&b'''=+0.0223
\end{array}
\]
Habentur itaque quatuor aequationes
\[
\begin{gathered}
+62''48-0.1744\mathrm d\AscendingNode+1.1957\mathrm di=0\\
+9.49+0.3578\mathrm d\AscendingNode-0.8172\mathrm di=0\\
-46.79-0.3292\mathrm d\AscendingNode-0.8685\mathrm di=0\\
-25.55-0.6915\mathrm d\AscendingNode+0.2034\mathrm di=0
\end{gathered}
\]
e quibus per principium quadratorum minimorum deducitur
\[
\begin{aligned}
\mathrm d\AscendingNode&=-54''41\\
\mathrm di&=-42.06
\end{aligned}
\]
ita ut habeatur
\begin{center}\fontsize{9}{11}\selectfont\setlength{\tabcolsep}{3pt}\begin{tabular}{lr}
longitudo nodi ascendentis pro initio anni 1803,&$172^0\,27'\,52''39$\\
inclinatio orbitae $\ .\ .\ .\ .\ .\ .\ .\ .\ .\ .$&$34\quad36\quad49.44$
\end{tabular}\end{center}

Elementa reliqua vel per comparationem eorum, quae in singulis tribus hypothesibus prodierunt, erui, vel quod accuratius est, per calculum novum longitudinum in orbita atque repetitionem operationum in art. 5 explicatarum determinari poterunt. Methodus prima suppeditat
\begin{center}\fontsize{9}{11}\selectfont\setlength{\tabcolsep}{3pt}\begin{tabular}{lr}
longitudinem perihelii 1803 $\ .\ .\ .\ .\ .\ .\ .$&$120^0\,58'\,3''86$\\
angulum $\varphi\ .\ .\ .\ .\ .\ .\ .\ .\ .\ .\ .\ .\ .$&$14\quad9\quad36.25$\\
longitudinem mediam in oppositione prima $\ .\ .$&$89\quad20\quad32.08$\\
motum medium diurnum sidereum $\ .\ .\ .\ .\ .\ .\ .$&$770''7899$\\
logarithmum semiaxis maioris. $\ .\ .\ .\ .\ ,\ .\ .$&$0.4420471$
\end{tabular}\end{center}
\SourceLeaf{20}
Per methodum alteram invenitur
\begin{center}\fontsize{9}{11}\selectfont\setlength{\tabcolsep}{3pt}\begin{tabular}{lr}
longitudo perihelii 1803 $\ .\ .\ .\ .\ .\ .\ .\ .$&$120^0\,58'\,4''81$\\
angulus $\varphi\ .\ .\ .\ .\ .\ .\ .\ .\ .\ .\ .\ .\ .$&$14\quad9\quad36.63$\\
longitudo media in oppositione prima $\ .\ .\ .\ .$&$89\quad20\quad31.81$\\
motus medius diurnus sidereus $\ .\ .\ .\ .\ .\ .\ .$&$770''7893$\\
logarithmus semiaxis maioris $\ .\ .\ .\ .\ .\ .\ .\ .$&$0.4420473$
\end{tabular}\end{center}
Cum his elementis conveniunt ea, quae supra (art. 3) tradidi.
\Article{9}
Quantumvis magnae sint perturbationes, quas Pallas a reliquis planetis patitur, tamen, experientia teste, elementa elliptica quatuor oppositionibus adaptata planetae motui intra integrum hoc tempus satis bene satisfaciunt: quin adeo etiam a motu antecedente ac consequente, nisi intervallum temporis nimium assumatur, parum differunt, ita ut e. g. elementa secunda in art. 3 tradita, in oppositione anni 1803 tribus, in oppositione anni 1809 quinque minutis primis a longitudine heliocentrica observata discrepaverint. Quae quum ita sint, ad construendam ephemeridem pro motu planetae futuro semper commodissimum mihi videtur elementis pure ellipticis uti, quae e quatuor oppositionibus proxime antecedentibus derivata erant, siquidem multitudo aequationum a perturbationibus oriundarum tanta certo evasura est, ut computus vel unius loci heliocentrici planetae tantum non aeque operosus evadere debeat, ac calculus elementorum ellipticorum per praecepta supra tradita. Ita locum planetae geocentricum salvo errore paucorum minutorum primorum tuto semper praedicere licebit, quae praecisio ad inveniendum planetam sufficit.
\Article{10}
Attamen postulat scientiae dignitas, ut consensui stabiliori prospiciatur, quem antequam perturbationes in calculum introducantur obtineri non posse manifestum est. Praematurus fuisset, meo quidem iudicio, calculus tam prolixus taediique plenus, quamdiu observationum copia tempus nimis exiguum complectebatur, perturbationesque a planetis reliquis oriundae vix se manifestabant. Nunc vero, ubi motus ellipticus non amplius sufficit, ad omnia loca observata inter se concilianda, tempus adesse videtur, ubi de theoria accuratiori cogitari potest. Quo pacto calculum perturbationum, quas Pallas patitur praesertim a Iove,
\SourceLeaf{21}
commodissime atque exactissime absolvendum censeam, quum methodis pro aliis planetis adhibitis propter nimiam excentricitatem atque inclinationem vix ac ne vix quidem uti liceat, mox alio loco fusius explicabo: de elementis autem ellipticis, quae maxime idonea videantur, ut calculus perturbationum ipsis superstruatur, in sequentibus adhuc agam. Eruenda scilicet mihi propono elementa elliptica, quae non his vel illis oppositionibus exacte, sed omnibus, quae hactenus observatae sunt, quam proxime satisfaciant. Methodum quidem, per quam tale negotium exsequi liceat, iam in Theoria Motus Corporum Coelestium art. 187 succincte descripsi: sed quum non solum ea, quae illic generaliter tractavi, in casu speciali, ubi loca observata sunt oppositiones, quaedam compendia admittant, sed etiam quaedam artificia practica, per quae applicationem methodi quadratorum minimorum faciliorem reddere iamdudum solitus sum, in opere illo desiderentur, astronomis haud ingratum fore spero, si hosce calculos aliquanto fusius hic tradidero. Quum totum negotium versetur in determinatione \textit{correctionum} elementis approximatis, quae ab omnibus locis observatis haud multum aberrare supponuntur, adiiciendarum, totus labor a duobus momentis pendebit: primo scilicet formari debent aequationes lineares, quas singula loca observata suppeditant, dein ex his aequationibus valores incognitarum maxime idonei sunt eruendi.
\Article{11}
Sit secundum elementa approximata
\begin{description}\setlength{\itemsep}{0pt}\setlength{\parsep}{0pt}
\item[$L$] longitudo media planetae pro epocha arbitraria
\item[$t$] numerus dierum inde ab epocha usque ad momentum observationis elapsorum
\item[$\MeanMotionGlyph$] motus medius diurnus sidereus in minutis secundis
\item[$\Pi$] longitudo perihelii
\item[$e=\sin\varphi$] excentricitas
\item[$a$] semiaxis maior
\item[$r$] radius vector
\item[$v$] anomalia vera
\item[$E$] anomalia excentrica
\item[$\AscendingNode$] longitudo nodi ascendentis
\item[$i$] inclinatio orbitae
\item[$u$] argumentum latitudinis
\item[$\lambda$] longitudo heliocentrica
